%%%%%%%%%%%%%%%%%%%%%%%%%%%%%%%%%%%%%%%%%%%%%%%%%%%%%%%%%%%%%%%%
\firstsection{引言}
\maketitle
\thispagestyle{plain}
%%%%%%%%%%%%%%%%%%%%%%%%%%%%%%%%%%%%%%%%%%%%%%%%%%%%%%%%%%%%%%%%
\subsection{研究背景与研究意义}

大语言模型正在以前所未有的速度渗透至经济决策的核心地带。
从自动化投资顾问和算法交易系统，到信贷审批与保险定价，LLM已从纯粹的信息处理工具演化为具有实质性决策权限的自主代理。
Lopez-Lira与Tang的先驱性工作表明，ChatGPT能够显著预测股票收益，其表现优于传统情感分析方法~\cite{LopezLira2023Can}。
Fatouros等人提出的MarketSenseAI框架利用GPT-4的链式推理能力，在S\&P 100股票上实现了10\%至30\%的超额收益~\cite{Fatouros2024Can}。
Yu等人设计的FinMem框架进一步通过分层记忆与角色设定构建了具备交易策略能力的LLM代理~\cite{Yu2023FinMem}。
Saha等人的综述系统梳理了LLM代理在投资管理中的应用版图，同时指出了鲁棒性、可解释性和实际部署等尚未解决的挑战~\cite{Saha2025Large}。
上述进展共同描绘了一幅清晰的图景，即LLM代理正在成为金融市场中不可忽视的参与力量。
与传统算法不同，LLM代理不仅输出数值预测，还会根据自然语言中的角色、责任和情境线索调整回答。
这使其风险偏好不再只是模型准确率或收益率问题，而成为一种可被实验操纵、也可能被提示词意外触发的行为特征。
因此，在部署前对LLM代理进行类似“行为审计”的实验测量，具有明确的现实必要性。

然而，这一趋势也催生了一个紧迫但尚未得到充分回答的问题——当LLM以代理人身份替他人做出风险决策时，其风险偏好是否会发生系统性偏移？
这一问题的重要性体现在三个层面。
第一，在实践层面，如果LLM在代理角色下表现出系统性保守或系统性激进的偏差，将直接影响数以亿计的投资者的资产安全。
第二，在理论层面，行为经济学已通过大量实验证实人类在替他人决策时存在显著的风险偏移效应，Polman与Wu的元分析纳入了71篇论文中的128个效应量（总计14443个观测值），确认了该效应的存在同时揭示了研究间的高度异质性~\cite{Polman2020Decision}。
然而，LLM作为一类缺乏真实财务痛感和声誉焦虑的新型决策主体，其代理决策行为是否遵循人类的行为规律，是一个有待检验的理论问题。
第三，在AI安全治理层面，主流LLM均经历了大规模的安全对齐训练。Ouyang等人通过人类反馈的强化学习（RLHF）使模型与用户意图对齐~\cite{Ouyang2022Training}，Bai等人提出的Constitutional AI方法进一步通过AI自我反馈实现无害化训练~\cite{Bai2022Constitutional}。
这些机制可能在代理场景中被异常激活，产生难以预见的风险偏好偏移。

为应对上述挑战，本文以Kahneman与Tversky的前景理论~\cite{Kahneman1979Prospect}和LLM行为偏差的最新发现为理论起点，对照相关文献，聚焦LLM安全对齐机制的潜在行为后果，通过严格的实验设计系统检验LLM在代理决策中的风险偏好变化。

\subsection{研究问题}

本研究聚焦的核心问题如下：

RQ1: 大语言模型在为他人代理风险决策时，相较于为自身决策，其风险偏好是否发生系统性偏移？

RQ2: 该偏移的方向与幅度是否受到损益情境和问责压力（委托人是否可观察决策过程）的调节？

RQ3: 不同模型、推理模式和题目呈现方式是否会改变上述效应的大小与稳定性？

围绕这些问题，本文的边际贡献主要体现在三方面。
第一，本文将人类“替他人决策”文献中的代理角色和问责线索迁移到LLM代理场景中，检验非人类决策主体是否呈现类似或相反的风险偏移。
第二，本文将Holt-Laury风险量表与损益域操纵结合，使风险偏好可以被转化为可比较的安全选择数量和转换点。
第三，本文同时比较多款模型，并加入深度思考模式和逐题独立调用检验，以避免将单一模型、单一Prompt下的表现误读为LLM的一般人格。

\subsection{理论逻辑与研究假设}

本研究的理论出发点如下。

第一点是 LLM 底层的安全对齐机制。
以RLHF和Constitutional AI为代表的训练范式，其核心目标之一是抑制模型生成有害、极端或高风险的建议~\cite{Ouyang2022Training,Bai2022Constitutional}。
该目标在预训练和微调阶段以参数级约束的形式嵌入模型权重。
当LLM被赋予代理角色，需要为委托人的利益做出风险判断时，安全对齐所塑造的保守倾向可能被角色设定中的责任暗示所放大——模型可能将替他人决策解读为一种需要格外审慎的情境，从而产生过度的风险规避。
我们将这一机制概括为算法谨慎偏移假设：
\begin{quote}
\textbf{H1A（算法谨慎偏移假设）：} 在收益域中，LLM以代理人身份替他人决策时，相较于为自己决策，将表现出更高的风险规避水平，表现为安全选择数量的显著增加。
\end{quote}

然而，这一假设与人类社会的相关研究不符，因此我们进一步进行探究。
行为经济学的实证研究反复表明，人类代理人在替他人决策时，其风险偏好的偏移方向取决于制度约束与心理机制的复杂交互。
Fullbrunn与Luhan发现有限责任制度（决策者仅参与收益而不承担损失）会压倒责任感，导致代理人为他人承担更高风险~\cite{Fullbrunn2019Responsibility}。
Stone等人的研究表明后悔顾虑在两种决策角色下均增加风险规避~\cite{Stone2002Risk}，而Batteux等人则发现代理决策中的风险偏好更接近期望值中性，主观风险偏好被削弱~\cite{Batteux2017Risk}。
对于LLM而言，它没有真实的财富损益体验，也不存在人类代理人的声誉焦虑或法律问责恐惧——至少在无观察者的条件下如此。
这意味着LLM在代理角色下可能出现与人类相反的偏移方向：由于缺乏对委托人损失的真切痛感，LLM可能对高风险选项中的高收益端赋予不成比例的注意力，从而表现出更强的风险寻求倾向。
我们将这一竞争性机制概括为道德风险偏移假设：
\begin{quote}
\textbf{H1B（道德风险偏移假设）：} 在收益域中，LLM以代理人身份替他人决策时，相较于为自己决策，将表现出更高的风险寻求水平，表现为安全选择数量的显著减少。
\end{quote}

H1A和H1B构成了一对相互排斥的核心假设，其实证检验结果将揭示安全对齐机制与道德风险机制在LLM代理决策中的相对主导地位。

第二点是问责压力的调节效应。
在人类代理决策中，当委托人能够观察代理人的每一步选择时，问责压力会显著改变代理人的行为模式。
Fullbrunn与Luhan在后续研究中进一步确认了责任感的调节作用~\cite{Fullbrunn2019Responsibility}，Buckle等人的实验则发现人均捐赠金额而非委托人数量产生了显著的谨慎效应~\cite{Buckle2024Needs}。
从理论逻辑推演，问责压力（可观察性）应当强化LLM在代理角色下的现有行为倾向：如果H1A成立（算法谨慎偏移主导），则可观察性应进一步放大谨慎程度；如果H1B成立（道德风险偏移主导），则可观察性应抑制过度的风险寻求。
然而，前景理论的一个经典预测是在损失域中，人类决策者的风险偏好会发生反转——在损失域中倾向于风险寻求，即所谓反射效应。
当代理角色与损失情境叠加时，问责压力的作用方向可能变得更为复杂：代理人在损失域中可能因为害怕被委托人责备而采取更为保守的策略以规避确定性损失，也可能因为试图翻盘而冒更大的风险。
据此，我们提出问责压力与损益域交互的调节假设：
\begin{quote}
\textbf{H2（问责压力与损益域交互调节假设）：} 问责压力（可观察性）在收益域与损失域中对LLM代理风险偏好的调节方向相反。在收益域中，可观察性将强化主导偏移方向（加剧H1A或抑制H1B）；在损失域中，可观察性将促使LLM向更加审慎的方向偏移，以减少对委托人的可见损失。
\end{quote}

第三点是模型间的异质性。
不同LLM在架构设计、预训练语料分布、RLHF实施方案和安全对齐强度上存在显著差异。
Qi等人的研究揭示即使是少量对抗性微调样本也能破坏LLM的安全对齐，甚至良性数据的微调也可能无意中降低安全性~\cite{Qi2024Finetuning}。
McIntosh等人则发现RLHF无法完全缓解意识形态条件化提示的影响，且其代表性不足以涵盖多元人类价值观~\cite{McIntosh2024Inadequacy}。
Liu等人在ACL上的工作表明经过RLHF微调的模型在人格引导生成中更易被操控~\cite{Liu2024Evaluating}。
Wu等人揭示了指令微调如何将预训练知识向用户导向任务旋转~\cite{Wu2024Language}。
这些发现共同指向一个可能性：不同厂商和不同训练策略下产出的LLM，其代理决策风险偏好可能存在系统性差异。
据此，我们提出模型间异质性假设：
\begin{quote}
\textbf{H3（模型间异质性假设）：} 不同架构和训练策略的LLM在代理决策风险偏好上存在显著差异，表现为模型间在相同实验条件下的安全选择数量和转换点分布上的统计显著性差别。
\end{quote}
